\section{Visual Elements}\label{sec:visual-elements}

\subsection{Figures and Images}\label{subsec:figures}
LaTeX can easily incorporate images using the \texttt{graphicx} package.

\subsubsection{Basic Figure Insertion}\label{subsubsec:basic-figures}
The basic syntax for including figures is:

\begin{verbatim}
\begin{figure}[<placement>]
    \centering
    \includegraphics[width=<width>]{<filename>}
    \caption{<caption text>}
    \label{fig:<label>}
\end{figure}
\end{verbatim}

Example:
\begin{figure}[htbp]
    \centering
    \includegraphics[width=0.4\textwidth]{images/default}
    \caption{A simple figure example}
    \label{fig:simple-example}
\end{figure}

\subsubsection{Subfigures}\label{subsubsec:subfigures}
The \texttt{subcaption} package allows multiple images in a single figure:

\begin{figure}[htbp]
    \centering
    \begin{subfigure}[b]{0.3\textwidth}
        \centering
        \includegraphics[width=\textwidth]{images/default}
        \caption{First subfigure}
        \label{fig:subfig-a}
    \end{subfigure}
    \hfill
    \begin{subfigure}[b]{0.3\textwidth}
        \centering
        \includegraphics[width=\textwidth]{images/default}
        \caption{Second subfigure}
        \label{fig:subfig-b}
    \end{subfigure}
    \caption{Multiple subfigures in one figure}
    \label{fig:subfigures}
\end{figure}

\subsubsection{Image Customization}\label{subsubsec:image-customization}
The \texttt{includegraphics} command accepts various options:
\begin{itemize}
    \item \texttt{width=0.8\textbackslash textwidth} - Set width (relative or absolute)
    \item \texttt{height=5cm} - Set height
    \item \texttt{scale=0.75} - Scale the image
    \item \texttt{angle=45} - Rotate the image
    \item \texttt{trim=left bottom right top} - Crop the image
    \item \texttt{clip} - Used with trim to enable cropping
\end{itemize}

\subsection{Tables}\label{subsec:tables}
Tables in LaTeX offer advanced formatting and alignment options.

\subsubsection{Basic Tables}\label{subsubsec:basic-tables}
The basic syntax for tables:

\begin{verbatim}
\begin{table}[<placement>]
    \caption{<caption text>}
    \centering
    \begin{tabular}{<column alignment>}
        \toprule
        <header row> \\
        \midrule
        <data rows> \\
        \bottomrule
    \end{tabular}
    \label{tab:<label>}
\end{table}
\end{verbatim}

Example:
\begin{table}[htbp]
    \caption{Sample Data Comparison}
    \centering
    \begin{tabular}{lcc}
        \toprule
        \textbf{Feature} & \textbf{Method A} & \textbf{Method B} \\
        \midrule
        Accuracy & 87.5\% & 92.3\% \\
        Speed & 1.2 s & 2.5 s \\
        Memory & 4.5 GB & 2.8 GB \\
        \bottomrule
    \end{tabular}
    \label{tab:sample-comparison}
\end{table}

\subsubsection{Column Alignment}\label{subsubsec:column-alignment}
Column specifiers:
\begin{itemize}
    \item \texttt{l} - Left-aligned column
    \item \texttt{c} - Center-aligned column
    \item \texttt{r} - Right-aligned column
    \item \texttt{p\{width\}} - Paragraph column with specified width
    \item \texttt{|} - Vertical line between columns (e.g., \texttt{|l|c|r|})
\end{itemize}

\subsubsection{Advanced Tables}\label{subsubsec:advanced-tables}
For more complex tables:

\begin{table}[htbp]
    \caption{Advanced Table Example}
    \centering
    \begin{tabular}{@{}l>{\raggedright\arraybackslash}p{3cm}c@{}}
        \toprule
        \textbf{Category} & \textbf{Description} & \textbf{Value} \\
        \midrule
        Type A & Short description that wraps to multiple lines & 42.5 \\
        \midrule
        Type B & Another description showing text wrapping & 18.3 \\
        \bottomrule
    \end{tabular}
    \label{tab:advanced-example}
\end{table}

\subsection{Color Usage}\label{subsec:color-usage}
The \texttt{xcolor} package provides comprehensive color handling.

\subsubsection{Text Color}\label{subsubsec:text-color}
Change text color using:
\begin{verbatim}
\textcolor{<color>}{<text>}
\end{verbatim}

Examples:
\textcolor{red}{Red text},
\textcolor{blue}{Blue text},
\textcolor{green}{Green text}

\subsubsection{Background Color}\label{subsubsec:background-color}
Add a background color with:
\begin{verbatim}
\colorbox{<color>}{<text>}
\end{verbatim}

Examples:
\colorbox{yellow}{Text with yellow background},
\colorbox{lightgray}{Text with gray background}

\subsubsection{Custom Colors}\label{subsubsec:custom-colors}
Define custom colors using RGB, HTML, or CMYK values:

\begin{verbatim}
\definecolor{customblue}{RGB}{0, 114, 122}
\definecolor{lightorange}{HTML}{FFE0B2}
\definecolor{mypurple}{cmyk}{0.3,0.8,0,0}
\end{verbatim}

Example: \textcolor{\bleuece}{This text uses the custom ECE blue color} defined in the preamble.

\subsubsection{Colored Boxes}\label{subsubsec:colored-boxes}
Create a colored box with text:

\begin{center}
\fcolorbox{black}{lightgray}{
    \begin{minipage}{0.8\textwidth}
        \centering
        This is a highlighted box that can contain important information\\
        or notices for the reader.
    \end{minipage}
}
\end{center}